\bookchapter{Chapitre 46 : Le chemin de la maîtrise — d'utilisateur à architecte}

\chapterintro{Une fois la méthodologie des chapitres 37 à 44 maîtrisée et les projets menés à bien avec l'IA, quelle est la suite ? Ce chapitre explore le passage d'« utilisateur d'outils IA » à « architecte de systèmes IA ». À ce stade, la question n'est plus « comment faire écrire du bon code à l'IA », mais « comment intégrer l'IA au mieux dans tout le processus de développement ».}

\booksection{Le protocole MCP et l'extension des outils}

MCP (Model Context Protocol) est un protocole ouvert apparu en 2025 et progressivement adopté en 2026 : il définit une norme de communication entre les modèles d'IA et les outils externes. En clair, MCP permet à l'IA d'« utiliser des outils » — pas seulement générer du code, mais aussi exécuter des commandes, interroger des bases de données, appeler des API et manipuler le système de fichiers.

Avant MCP, chaque outil IA avait sa propre façon d'intégrer des outils, incompatible avec les autres. L'arrivée de MCP change la donne : un outil conforme à MCP peut être utilisé par n'importe quel outil IA qui le prend en charge. Le coût d'intégration s'effondre, et le périmètre de l'IA s'étend de la « génération de texte » à l'« exécution d'actions ».

Pour l'utilisateur avancé, l'intérêt de MCP est là : tu peux doter l'IA d'outils sur mesure. Par exemple, si ton équipe possède un système de déploiement interne, tu peux écrire un outil MCP qui permette à l'IA d'appeler directement ce système, sans devoir passer manuellement par le terminal. Cette capacité de l'IA à « appeler des outils » est l'infrastructure de base de l'Agentic Engineering.

\booksection{Workflows personnalisés et Skills}

Quand tu répètes certains processus de programmation IA bien rodés, tu peux les figer en « workflows personnalisés » ou « Skills ». Par exemple, ton équipe a peut-être une procédure de relecture standard — d'abord la sécurité, puis la performance, puis la maintenabilité. Écris cette procédure sous forme de Skill, et l'IA la suivra automatiquement à chaque relecture de code.

La valeur centrale d'un workflow personnalisé, c'est la « cohérence » : il garantit que chaque exécution suit les mêmes critères, sans dérive liée à la longueur de la conversation ou au contexte. C'est particulièrement précieux en équipe : quand tous les membres utilisent le même workflow, le style et la qualité du code convergent.

La méthode varie selon l'outil. Cursor prend en charge les Rules et Custom Instructions, Claude Code les commandes personnalisées, Replit Runme. Mais l'idée reste la même : « coder » tes meilleures pratiques en modèles d'instructions que l'IA peut exécuter.

\booksection{Le mode de collaboration multi-agents}

En 2026, la collaboration multi-agents devient un sujet très discuté. Son idée-force : ne pas confier la tâche à une seule IA, mais faire collaborer plusieurs IA aux rôles distincts. Par exemple, une IA écrit le code, une autre écrit les tests, une autre relit — chacune son métier, et elles se contrôlent mutuellement.

Les atouts du multi-agents : la « spécialisation » et la « vérification croisée ». Chaque agent se concentre sur son domaine et produit de meilleure qualité ; les contrôles croisés entre agents détectent des problèmes qu'on pourrait laisser passer. Mais les limites sont nettes : la coordination coûte cher, la communication alourdit le travail et le débogage se complique.

\begin{tcbinfo}{Astuce : le « multi-agents » en solo, c'est possible}
En tant que développeur solo, pas besoin de monter un système multi-agents complexe — tu peux « jouer » différents rôles pour simuler la collaboration. Par exemple : une première conversation où l'IA écrit le code fonctionnel ; puis une nouvelle conversation (avec le contexte de CLAUDE.md) où l'IA « relit ce code en tant que relecteur » ; enfin une troisième où l'IA « corrige le code d'après les remarques de relecture ». Ce « one-man-show à plusieurs rôles » te donne l'essentiel des bénéfices du multi-agents, sans infrastructure sophistiquée.

\end{tcbinfo}

\booksection{Le chemin de l'apprentissage continu}

La programmation IA est un domaine qui évolue à grande vitesse. Les meilleures pratiques d'aujourd'hui seront peut-être à réviser demain. Quelques conseils pour continuer à apprendre :

Premier conseil : va aux sources plutôt qu'aux interprétations de seconde main. Les messages originaux d'Andrej Karpathy et son CLAUDE.md, la documentation officielle d'Anthropic et d'OpenAI, les notes de version des grands outils — voilà l'information de première main. Les articles et tutoriels de la communauté les interprètent et les résument : ils ont de la valeur, mais peuvent déformer. Prends l'habitude de lire les sources.

Deuxième conseil : apprends par la pratique. Ne fais pas « d'abord apprendre, ensuite faire », mais « apprendre en faisant ». Chaque nouveau type de problème (premier envoi de fichiers, première intégration d'une API tierce, par exemple) est une occasion d'acquérir les connaissances correspondantes. On apprend bien plus vite en situation qu'en théorie décontextualisée.

Troisième conseil : reste à l'affût sans courir après chaque nouveauté. Tous les un ou deux mois, consacre une demi-journée aux derniers outils et tendances — mais ne change pas d'outils à chaque sortie. Le critère : le nouvel outil lève-t-il vraiment l'un de tes goulots d'étranglement ? Sinon, ta combinaison actuelle suffit.

Quatrième conseil : construis ta base de connaissances personnelle. Note chaque conversation qui t'a aidé à résoudre un problème complexe, chaque leçon tirée d'un piège, chaque méthode qui a marché. Inutile de publier — ce sont tes « notes de méthode » à toi. Avec le temps, elles deviendront ton bien le plus précieux.

\begin{tcbpractice}{Passe à l'action : ouvre un « carnet de galères »}
Crée dès aujourd'hui un document (logiciel de notes ou simple fichier .md) et note les trois leçons de programmation IA que tu veux surtout retenir à ce jour — même un simple « réfléchir avant de demander ». Chaque galère, chaque nouvelle astuce : ajoute une ligne. Dans six mois, en le relisant, tu remercieras le toi d'aujourd'hui.

\end{tcbpractice}

\begin{tcbsummary}{L'essentiel du chapitre}
\begin{itemize}
  \item MCP est la norme de communication entre l'IA et les outils : elle permet d'exécuter des commandes, d'interroger des bases de données et d'appeler des API
  \item La valeur centrale des workflows/Skills : la cohérence — coder les meilleures pratiques en modèles d'instructions
  \item Multi-agents = spécialisation + vérification croisée ; en solo, on peut le simuler en jouant plusieurs rôles
  \item Quatre clés de l'apprentissage continu : lire les sources, apprendre en faisant, ne pas courir après la nouveauté, bâtir sa base de connaissances
\end{itemize}

\end{tcbsummary}
